% =============================================================================
% Statistical Test Finder — Decision Tree
% Author:  Tim Draws <drawstim@gmail.com>  |  https://timdraws.net
% License: MIT  (https://opensource.org/licenses/MIT)
% Source:  https://timdraws.net/files/StatisticalTestFinder.pdf
% =============================================================================
%
% ARCHITECTURE OVERVIEW
% ---------------------
%  1. Preamble        — packages, page geometry, hyperlink config
%  2. Forest styles   — reusable node/edge style definitions
%  3. Page style      — header/footer setup
%  4. Legend tree     — small forest that explains the column layout
%  5. Main tree       — the full decision tree (hypothesis tests only;
%                       descriptive and predictive analyses excluded by design)
%
% TREE COLUMN STRUCTURE (left → right)
%   Analysis Type → #DVs → DV type → #DV categories → #IVs
%   → IV type → #IV categories → Research Design
%   → Recommended Analysis → Non-Parametric Alternative
%
% KNOWN OMISSIONS (intentional, see Note 1 in footer)
%   - Mixed/random-effects models
%   - Polynomial & multinomial probit regression
%   - Most ML methods (not typically used for hypothesis testing)
% =============================================================================


% -----------------------------------------------------------------------------
% 1. PACKAGES & PAGE GEOMETRY
% -----------------------------------------------------------------------------

\documentclass{article}

\usepackage[utf8]{inputenc}

% Oversized "poster" page to accommodate the wide forest tree.
% heightrounded avoids geometry warnings about fractional lengths.
\usepackage[
  paperheight = 1850pt,
  paperwidth  = 1650pt,
  margin      = 50pt,
  heightrounded
]{geometry}

\geometry{
  bottom   = 150pt,
  footskip = 50pt
}

\usepackage{forest}      % Decision-tree drawing
\usepackage{fancyhdr}    % Custom header/footer
\usepackage{hyperref}    % Clickable URLs and cross-references


% -----------------------------------------------------------------------------
% 2. HYPERLINK CONFIGURATION
% -----------------------------------------------------------------------------

% We want underlined links (not coloured) so the tree stays monochrome.
% The two \Hy@AtBeginDocument overrides are necessary because hyperref's
% colorlinks=false still draws a coloured box by default; the patch below
% replaces the box border with a thin underline.
\makeatletter
\Hy@AtBeginDocument{%
  \def\@pdfborder{0 0 1}%           Border dimensions: x y w
  \def\@pdfborderstyle{/S/U/W 1}%   /S/U = Underline style, /W 1 = 1pt wide
}
\makeatother

\hypersetup{
  colorlinks     = false,   % Colour handled via border patch above
  linkcolor      = black,   % Fallback link text colour (rarely visible)
  linkbordercolor = black,  % Border colour for internal links
}


% -----------------------------------------------------------------------------
% 3. FOREST NODE/EDGE STYLE DEFINITIONS
% -----------------------------------------------------------------------------

\forestset{%
  % Internal (decision) nodes — draws a visible box.
  Node/.style={draw=black},%
  % Legacy alias used in the legend tree.
  L2/.style={edge={line width=0.8pt}},%
  % Leaf: recommended-analysis column — fixed width for alignment.
  Leaf/.style={draw=black, minimum width=60mm},%
  % LeafAnalysis: legend column headers — same width, no border.
  LeafAnalysis/.style={minimum width=60mm},%
}


% -----------------------------------------------------------------------------
% 4. PAGE STYLE (header suppressed; footer holds footnotes + credits)
% -----------------------------------------------------------------------------

\title{Statistical Test Finder}
\author{}
\date{}

\pagestyle{fancy}
\fancyhead{}                        % No header content
\renewcommand{\headrulewidth}{0pt}  % Remove header rule

% Footer: a single wide tabular cell that spans the full page width and
% holds pseudo-footnotes (numbered manually via \footnotemark) plus
% general notes and the licence statement.

\cfoot{%
  \begin{tabular}{p{2000pt}}
    \\
    % Footnote 1 — likelihood-ratio test siblings
    \footnotemark[1]{%
      Two closely related tests (approximations of the likelihood ratio
      test) are the \underline{\href{https://en.wikipedia.org/wiki/Wald_test}{Wald
      test}} and the \underline{\href{https://en.wikipedia.org/wiki/Score_test}{score test}}.
    }\\
    % Footnote 2 — one-sample t vs Z distinction
    \footnotemark[2]{%
      When testing whether a population mean equals a given value, use the
      one-sample $t$-test when the true population variance is unknown and the $Z$-test when it is known.
    }\\
    % Footnote 3 — Mann-Whitney alias
    \footnotemark[3]{%
      The Mann-Whitney $U$-test is also known as the Wilcoxon rank-sum
      test.
    }\\
    % Footnote 4 — Kruskal-Wallis limitation
    \footnotemark[4]{%
      The Kruskal-Wallis test can only test main effects (no interaction
      effects).
    }\\\\
    % --- General notes ---
    \textbf{Note 1.} This sheet does not cover every single existing statistical hypothesis test. To avoid overload, infrequently used methods (e.g., polynomial regression or multinomial probit regression),\\ \hspace{0.5cm} mixed/random-effect models, and predictive methods not typically used for hypothesis testing (e.g., tree-based ML) are intentionally omitted.\\
    \textbf{Note 2.} All tests are presented as classical (frequentist)
    hypothesis tests. However, note that many of them have Bayesian counterparts implemented in R, Python, or JASP.\\\\
    % --- Licence & provenance ---
    Originally created by \underline{\href{https://timdraws.net}{Tim Draws}}. Released under the \underline{\href{https://opensource.org/licenses/MIT}{MIT License}}. Last update: \today. Links: \underline{\href{https://timdraws.net/files/StatisticalTestFinder.pdf}{latest version}}, \underline{\href{https://github.com/timdraws/StatisticalTestFinder}{GitHub repository}} -- feedback/changes/additions welcome!
  \end{tabular}%
}

\pagenumbering{gobble}  % Suppress page numbers on this poster-style sheet


% =============================================================================
% DOCUMENT BODY
% =============================================================================

\begin{document}


% -----------------------------------------------------------------------------
% 5. LEGEND TREE
% Renders a small "key" showing what each column in the main tree represents.
% Uses the same forest settings as the main tree so column widths align when
% both forests are placed on the same page.
% -----------------------------------------------------------------------------

\begin{forest}
  for tree={%
    grow=0, reversed,% left-to-right layout
    parent anchor=east, child anchor=west,%
    edge={line cap=round}, outer sep=+1pt,%
    rounded corners, minimum width=40mm, minimum height=30mm,%
    l sep=10mm% horizontal level spacing
  }
  [\textbf{(Start Here)}\\Analysis Type, align=center, base=center, L2
    [Number of\\Dependent Variables, align=center, base=center, L2
      [Type of\\Dependent Variable(s), align=center, base=center, L2
        [(If Categorical DV)\\Number of Categories, align=center, base=center
          [Number of\\Independent Variables, align=center, base=center, L2
            [Type of\\Independent Variable(s), align=center, base=center, L2
              [(If Categorical IV)\\Number of Categories, align=center, base=center
                [(If Categorical IV)\\Research Design, align=center, base=center
                  [\textbf{Recommended Analysis}, align=center, base=center, LeafAnalysis
                    [\textbf{Non-Parametric Alternative}, align=center, base=center, LeafAnalysis]
                  ]
                ]
              ]
            ]
          ]
        ]
      ]
    ]
  ]
\end{forest}


% -----------------------------------------------------------------------------
% 6. MAIN DECISION TREE
%
% Structure mirrors the legend columns above. Tier names are used to keep
% sibling nodes vertically aligned across branches:
%   AnalysisType → nDV → DVtype → nCatDV → nIV → IVtype → nCat
%   → resDesign → recAnalysis → nonPara
%
% STYLE CONVENTIONS
%   Node       — all branch (decision) nodes
%   Leaf       — recommended-analysis leaf (rightmost substantive column)
%   nonPara    — non-parametric alternative leaf (child of Leaf)
%   tier       — alignment anchor; must match the legend column position
%
% FOOTNOTE SUPERSCRIPTS
%   $^1$ Likelihood-ratio test siblings (Wald, Score)
%   $^2$ t-test vs Z-test distinction
%   $^3$ Mann-Whitney = Wilcoxon rank-sum alias
%   $^4$ Kruskal-Wallis: main effects only
% -----------------------------------------------------------------------------

\begin{forest}
  for tree={%
    grow=0, reversed,%
    parent anchor=east, child anchor=west,%
    edge={line cap=round}, outer sep=+1pt,%
    rounded corners, minimum width=40mm, minimum height=8mm,%
    l sep=10mm%
  }
  % ============================================================
  % ROOT
  % ============================================================
  [Hypothesis Test, tier=AnalysisType, Node
    % ----------------------------------------------------------
    % BRANCH A: ONE dependent variable
    % ----------------------------------------------------------
    [One, tier=nDV, Node
      % ··· A-1: NOMINAL dependent variable ···················
      [Nominal, tier=DVtype, Node
        % A-1a: Binary DV (2 categories)
        [Two, tier=nCatDV, Node
          % No IV — test against a fixed probability
          [None (Test Value), tier=nIV, Node
            [\href{http://sites.utexas.edu/sos/guided/inferential/categorical/univariate/binomial/}{Binomial Test},
              tier=recAnalysis, Leaf]
          ]
          % One IV
          [One, tier=nIV, Node
            % IV is nominal
            [Nominal, tier=IVtype, Node
              % Binary nominal IV
              [Two, tier=nCat, Node
                [Between-Subjects, tier=resDesign, Node
                  [\href{https://en.wikipedia.org/wiki/Chi-squared_test}{Fisher's Exact Test},
                    tier=recAnalysis, Leaf]
                ]
                [Within-Subjects, tier=resDesign, Node
                  [\href{https://www.statstest.com/mcnemar-test/}{McNemar's Test},
                    tier=recAnalysis, Leaf]
                ]
              ]
              % Nominal IV (2+ categories)
              [Two or more, tier=nCat, Node
                [\href{https://en.wikipedia.org/wiki/Chi-squared_test}{Chi-Squared Test},
                  tier=recAnalysis, Leaf]
                [\href{https://en.wikipedia.org/wiki/Likelihood-ratio_test}{Likelihood Ratio Test}$^1$,
                  tier=recAnalysis, Leaf]
              ]
            ] % end Nominal IV type
            % IV is ordinal
            [Ordinal, tier=IVtype, Node
              [\href{https://www.statisticshowto.com/rank-biserial-correlation/}{Rank-Biserial Correlation},
                tier=recAnalysis, Leaf]
            ]
            % IV is continuous
            [Continuous, tier=IVtype, Node
              [\href{https://en.wikipedia.org/wiki/Logistic_regression}{Logistic Regression},
                tier=recAnalysis, Leaf]
              [\href{https://www.andrews.edu/~calkins/math/edrm611/edrm13.htm}{Point-Biserial Correlation},
                tier=recAnalysis, Leaf]
            ]
          ] % end One IV (binary DV)
          % Two or more IVs
          [Two or more, tier=nIV, Node
            [Continuous, tier=IVtype, Node
              [\href{https://www.aphis.usda.gov/animal_health/surveillance_toolbox/docs/epi_surv/discriminant_analysis_statsoft_textbook.pdf}{Discriminant Analysis},
                tier=recAnalysis, Leaf]
            ]
            [Categorical, tier=IVtype, Node
              [\href{https://en.wikipedia.org/wiki/Log-linear_analysis}{Loglinear Analysis},
                tier=recAnalysis, Leaf]
            ]
            [Any, tier=IVtype, Node
              [\href{https://en.wikipedia.org/wiki/Logistic_regression}{Logistic Regression},
                tier=recAnalysis, Leaf]
            ]
          ] % end Two or more IVs (binary DV)
        ] % end A-1a (binary DV)
        % A-1b: Polytomous DV (3+ categories)
        [More than two, tier=nCatDV, Node
          % No IV — test against expected distribution
          [None (Distribution), tier=nIV, Node
            [\href{https://www.statology.org/multinomial-test/}{Multinomial Test},
              tier=recAnalysis, Leaf]
          ]
          % One IV
          [One, tier=nIV, Node
            [Categorical, tier=IVtype, Node
              [\href{https://en.wikipedia.org/wiki/Chi-squared_test}{Chi-Squared Test},
                tier=recAnalysis, Leaf]
              [\href{https://en.wikipedia.org/wiki/Likelihood-ratio_test}{Likelihood Ratio Test}$^1$,
                tier=recAnalysis, Leaf]
            ]
            [Any, tier=IVtype, Node
              [\href{https://stats.idre.ucla.edu/r/dae/multinomial-logistic-regression/}{Multinomial Logistic Regression},
                tier=recAnalysis, Leaf]
            ]
          ]
          % Two or more IVs
          [Two or more, tier=nIV, Node
            [Continuous, tier=IVtype, Node
              [\href{https://www.aphis.usda.gov/animal_health/surveillance_toolbox/docs/epi_surv/discriminant_analysis_statsoft_textbook.pdf}{Discriminant Analysis},
                tier=recAnalysis, Leaf]
            ]
            [Categorical, tier=IVtype, Node
              [\href{https://en.wikipedia.org/wiki/Log-linear_analysis}{Loglinear Analysis},
                tier=recAnalysis, Leaf]
            ]
            [Any, tier=IVtype, Node
              [\href{https://stats.idre.ucla.edu/r/dae/multinomial-logistic-regression/}{Multinomial Logistic Regression},
                tier=recAnalysis, Leaf]
            ]
          ]
        ] % end A-1b (polytomous DV)
      ] % end A-1 (Nominal DV)
      % ··· A-2: ORDINAL dependent variable ···················
      [Ordinal, tier=DVtype, Node
        % One IV
        [One, tier=nIV, Node
          [Nominal, tier=IVtype, Node
            % NOTE: only binary nominal IV shown; polytomous case → Kruskal-Wallis
            % (covered under Continuous DV branch)
            [Two, tier=nCat, Node
              [\href{https://www.statisticshowto.com/rank-biserial-correlation/}{Rank-Biserial Correlation},
                tier=recAnalysis, Leaf]
              [\href{https://www.statology.org/mann-whitney-u-test}{Mann-Whitney $U$-Test$^3$},
                tier=recAnalysis, Leaf]
            ]
          ]
          [Ordinal, tier=IVtype, Node
            [\href{https://en.wikipedia.org/wiki/Kendall_rank_correlation_coefficient}{Kendall's $\tau$},
              tier=recAnalysis, Leaf]
            [\href{https://en.wikipedia.org/wiki/Spearman's_rank_correlation_coefficient}{Spearman's $\rho$},
              tier=recAnalysis, Leaf]
          ]
          [Continuous, tier=IVtype, Node
            [\href{https://en.wikipedia.org/wiki/Ordered_logit}{Ordinal Logistic Regression},
              tier=recAnalysis, Leaf]
            [\href{https://www.andrews.edu/~calkins/math/edrm611/edrm13.htm}{Biserial Correlation},
              tier=recAnalysis, Leaf]
          ]
        ] % end One IV (ordinal DV)
        % Two or more IVs
        [Two or more, tier=nIV, Node
          [\href{https://en.wikipedia.org/wiki/Ordered_logit}{Ordinal Logistic Regression},
            tier=recAnalysis, Leaf]
        ]
      ] % end A-2 (Ordinal DV)
      % ··· A-3: CONTINUOUS dependent variable ················
      [Continuous, tier=DVtype, Node
        % No IV — one-sample tests
        [None (Test Value), tier=nIV, Node
          [\href{https://www.statology.org/one-sample-t-test/}{One Sample $t$-Test},
            tier=recAnalysis, Leaf
            % Non-parametric alternative
            [\href{https://en.wikipedia.org/wiki/Wilcoxon_signed-rank_test}{Wilcoxon Signed-Rank Test},
              tier=nonPara, Leaf]
          ]
          % Z-test: only when true population variance is known (footnote 2)
          [\href{https://courses.washington.edu/psy315/tutorials/z_test_tutorial.pdf}{$Z$-Test$^2$},
            tier=recAnalysis, Leaf]
        ]
        % One IV
        [One, tier=nIV, Node
          % IV is categorical
          [Categorical, tier=IVtype, Node
            % Binary categorical IV
            [Two, tier=nCat, Node
              [Between-Subjects, tier=resDesign, Node
                [\href{https://www.statstest.com/independent-samples-t-test/}{Independent Samples $t$-Test},
                  tier=recAnalysis, Leaf
                  [\href{https://www.statology.org/mann-whitney-u-test/}{Mann-Whitney $U$-Test$^3$},
                    tier=nonPara, Leaf]
                ]
                % Welch's t-test: preferable when variances are unequal
                [\href{https://www.statology.org/welchs-t-test/}{Welch's $t$-Test},
                  tier=recAnalysis, Leaf]
                [\href{https://www.andrews.edu/~calkins/math/edrm611/edrm13.htm}{Point-Biserial Correlation},
                  tier=recAnalysis, Leaf]
              ]
              [Within-Subjects, tier=resDesign, Node
                [\href{https://www.statology.org/paired-samples-t-test/}{Paired Samples $t$-Test},
                  tier=recAnalysis, Leaf
                  [\href{https://www.statstest.com/wilcoxon-signed-rank-test/}{Wilcoxon Signed-Rank Test},
                    tier=nonPara, Leaf]
                ]
              ]
            ] % end binary categorical IV
            % Polytomous categorical IV (3+ groups)
            [More Than Two, tier=nCat, Node
              [Between-Subjects, tier=resDesign, Node
                [\href{https://www.statology.org/one-way-anova/}{One-Way ANOVA},
                  tier=recAnalysis, Leaf
                  [\href{https://www.statology.org/kruskal-wallis-test/}{Kruskal-Wallis Test},
                    tier=nonPara, Leaf]
                ]
              ]
              [Within-Subjects, tier=resDesign, Node
                [\href{https://www.statology.org/repeated-measures-anova/}{One-Way Repeated Measures ANOVA},
                  tier=recAnalysis, Leaf
                  [\href{https://www.statology.org/friedman-test/}{Friedman Test},
                    tier=nonPara, Leaf]
                ]
              ]
            ] % end polytomous categorical IV
          ] % end Categorical IV type (one IV, continuous DV)
          % IV is continuous
          [Continuous, tier=IVtype, Node
            [\href{https://statistics.laerd.com/statistical-guides/pearson-correlation-coefficient-statistical-guide.php}{Pearson Correlation},
              tier=recAnalysis, Leaf
              [\href{https://en.wikipedia.org/wiki/Kendall_rank_correlation_coefficient}{Kendall's $\tau$},
                tier=nonPara, Leaf]
              [\href{https://en.wikipedia.org/wiki/Spearman's_rank_correlation_coefficient}{Spearman's $\rho$},
                tier=nonPara, Leaf]
            ]
            [\href{http://www.stat.yale.edu/Courses/1997-98/101/linreg.htm}{Simple Linear Regression},
              tier=recAnalysis, Leaf
              [\href{https://rcompanion.org/handbook/F_12.html}{Nonparametric Regression},
                tier=nonPara, Leaf]
            ]
          ]
        ] % end One IV (continuous DV)
        % Two or more IVs
        [Two or more, tier=nIV, Node
          % All IVs categorical (factorial ANOVA family)
          [Categorical, tier=IVtype, Node
            [Between-Subjects, tier=resDesign, Node
              [\href{https://sites.ualberta.ca/~lkgray/uploads/7/3/6/2/7362679/slides_-_multi-way_anovaa.pdf}{Multi-Way ANOVA},
                tier=recAnalysis, Leaf
                % Kruskal-Wallis: main effects only (footnote 4)
                [\href{https://www.statology.org/kruskal-wallis-test/}{Kruskal-Wallis Test$^4$},
                  tier=nonPara, Leaf]
              ]
            ]
            [Within-Subjects, tier=resDesign, Node
              [\href{https://biostats.w.uib.no/factorial-repeated-measures-anova-two-way-repeated-measures-anova/}{Factorial Repeated Measures ANOVA},
                tier=recAnalysis, Leaf]
            ]
            [Mixed, tier=resDesign, Node
              [\href{http://web.pdx.edu/~newsomj/uvclass/ho_mixed.pdf}{Factorial Mixed ANOVA},
                tier=recAnalysis, Leaf]
            ]
          ] % end Categorical IVs
          % All IVs continuous (regression family)
          [Continuous, tier=IVtype, Node
            [\href{http://www.stat.yale.edu/Courses/1997-98/101/linmult.htm}{Multiple Linear Regression},
              tier=recAnalysis, Leaf
              [\href{https://rcompanion.org/handbook/F_12.html}{Nonparametric Regression},
                tier=nonPara, Leaf]
            ]
            [\href{https://www.statstest.com/partial-correlation/}{Partial Correlation},
              tier=recAnalysis, Leaf]
          ]
          % Mixed IVs: categorical + continuous (ANCOVA family)
          [Both, tier=IVtype, Node
            [\href{http://www.stat.yale.edu/Courses/1997-98/101/linmult.htm}{Multiple Linear Regression},
              tier=recAnalysis, Leaf
              [\href{https://rcompanion.org/handbook/F_12.html}{Nonparametric Regression},
                tier=nonPara, Leaf]
            ]
            [\href{https://www.statology.org/ancova/}{ANCOVA},
              tier=recAnalysis, Leaf
              [\href{https://www.youtube.com/watch?v=jozx59H5IHw}{Quade's ANCOVA},
                tier=nonPara, Leaf]
            ]
          ]
        ] % end Two or more IVs (continuous DV)
      ] % end A-3 (Continuous DV)
    ] % end BRANCH A (one DV)
    % ----------------------------------------------------------
    % BRANCH B: TWO OR MORE dependent variables (multivariate)
    % Note: categorical/mixed DV types omitted — no standard
    % single analysis covers them.
    % ----------------------------------------------------------
    [Two or more, tier=nDV, Node
      % All DVs continuous
      [Continuous, tier=DVtype, Node
        % One IV
        [One, tier=nIV, Node
          [Categorical, tier=IVtype, Node
            [\href{https://www.statology.org/manova-spss/}{MANOVA},
              tier=recAnalysis, Leaf
              [\href{https://link.springer.com/article/10.1007/s10463-019-00717-3}{Nonparametric MANOVA},
                tier=nonPara, Leaf]
            ]
          ]
        ]
        % Two or more IVs
        [Two or more, tier=nIV, Node
          [Continuous, tier=IVtype, Node
            [\href{https://online.stat.psu.edu/stat505/lesson/13}{Canonical Correlation Analysis},
              tier=recAnalysis, Leaf]
          ]
          [Categorical, tier=IVtype, Node
            [\href{https://core.ecu.edu/wuenschk/MV/MANOVA/MANOVA2.pdf}{Factorial MANOVA},
              tier=recAnalysis, Leaf]
          ]
          [Both, tier=IVtype, Node
            [\href{https://www.statisticssolutions.com/free-resources/directory-of-statistical-analyses/multivariate-analysis-of-covariance-mancova/}{MANCOVA},
              tier=recAnalysis, Leaf]
          ]
        ]
      ] % end Continuous DVs
    ] % end BRANCH B (two or more DVs)
  ] % end ROOT (Hypothesis Test)
\end{forest}

\end{document}